\section{Trabajos futuros}

Como líneas de trabajo futuro, se propone ampliar ParkLink mediante la incorporación de sensores IoT para detectar la ocupación de las cocheras, modelos de predicción de la demanda y mecanismos de integración con operadores de estacionamientos. Asimismo, se busca desarrollar una solución escalable que permita su implementación en ciudades densamente pobladas de Latinoamérica y de los Estados Unidos, considerando las particularidades de movilidad, infraestructura y regulación de cada contexto.

La incorporación de sensores permitiría obtener información en tiempo real sobre la disponibilidad de los espacios y mejorar la precisión de los datos mostrados a los usuarios. A partir de esta información, podrían desarrollarse modelos capaces de analizar patrones de uso, anticipar los momentos de mayor demanda y recomendar alternativas de estacionamiento de acuerdo con la ubicación, el horario y las preferencias de cada persona. También sería conveniente evaluar distintos mecanismos de actualización y comunicación para garantizar que la información se mantenga confiable y oportuna.

Otra línea de trabajo consiste en establecer acuerdos con operadores de estacionamientos, comercios y organismos públicos para ampliar la cobertura de la plataforma. Estas integraciones podrían facilitar la incorporación de nuevas cocheras, la gestión de reservas y la implementación de métodos de pago digitales. Además, permitirían explorar indicadores de desempeño relacionados con la disponibilidad, el tiempo de búsqueda, la rotación de los espacios y el impacto de la solución sobre la congestión vehicular.

Estas funcionalidades y objetivos de expansión no forman parte del MVP propuesto, dado que requieren infraestructura tecnológica, disponibilidad de datos históricos, análisis de viabilidad y procesos de validación adicionales que exceden el alcance inicial del proyecto.

En etapas posteriores también será necesario realizar pruebas piloto en diferentes zonas y comparar los resultados obtenidos en contextos urbanos diversos. La evaluación deberá contemplar aspectos técnicos, económicos y de experiencia de usuario, así como cuestiones vinculadas con la privacidad, la seguridad de la información y el cumplimiento de la normativa vigente. De este modo, ParkLink podrá evolucionar de manera gradual, priorizando las funcionalidades con mayor impacto y asegurando que su crecimiento sea sostenible y adaptable a las necesidades de cada ciudad.
